\section{S-Track Contributions}
\label{sec:track}

The S-track comprises four papers (S.1--S.4) that implement the
three-component framework with worked examples.

\begin{table}[h]
\centering
\caption{S-track paper summary. Method and primary contribution for each paper.}
\label{tab:track-overview}
\begin{tabular}{llll}
\toprule
Paper & Title & Method & Primary contribution \\
\midrule
S.1 & Hypothesis testing & Permutation test & Calibrated p-values from \\
    & for gerrymandering & + calibrated ensemble & G.1 ensemble data \\
\midrule
S.2 & Bayesian redistricting & Bayesian inference & Bayesian Detection Score \\
    &                        & Beta posterior & (BDS) for courts \\
\midrule
S.3 & Power analysis for & Power-ESS formula & Min.\ ensemble sizes; \\
    & ensemble sizing & (Section 3.2) & 50-state table \\
\midrule
S.4 & Multiple testing & Benjamini-Hochberg & BH applied to L-track \\
    & FDR control & FDR at $q = 0.10$ & 1,050-test battery \\
\bottomrule
\end{tabular}
\end{table}

\subsection{S.1: Hypothesis Testing}

S.1 develops a permutation test framework for partisan gerrymandering
that converts G.1's ensemble percentiles into calibrated p-values.
The key step is accounting for ensemble autocorrelation (C2) and
tail-estimation uncertainty (C4) simultaneously.

For North Carolina 2022 enacted plan:
\begin{itemize}
\item Nominal percentile (G.1): 0.4th percentile (Democratic seats)
\item ESS correction (G.4): $n^* = 695$ (not 10,000)
\item Calibrated p-value: $p_{\rm cal} = 0.008$
\item 95\% CI on true percentile: $[0.2\%, 1.5\%]$
\end{itemize}

The calibrated p-value of 0.008 supports rejection at $\alpha = 0.05$
even after correction, but the confidence interval shows the claim is
``at most 1.5th percentile'' rather than exactly ``0.4th percentile.''

\subsection{S.2: Bayesian Redistricting}

S.2 provides the Bayesian Detection Score (BDS) developed in the
S.2 paper~\citep{s2paper}.
BDS converts the ensemble percentile and ESS into a posterior probability
that the plan was intentionally selected for partisan advantage.
For NC 2022 ($\hat{p} = 0.004$, $n^* \approx 695$): BDS $\approx 0.97$.

\subsection{S.3: Power Analysis}

S.3 derives minimum ensemble sizes for 80\% power at $\alpha = 0.05$
to detect a 10pp partisan advantage.
Key results (from the formula in Section~\ref{sec:framework}):
\begin{itemize}
\item NC ($k = 14$, $\rho_1 = 0.87$): minimum $\approx 5{,}000$ steps
\item CA ($k = 52$, $\rho_1 = 0.94$): minimum $\approx 30{,}000$ steps
\end{itemize}
Current practice (10,000 steps) gives 80\% power for NC but only 45\%
for California.

\subsection{S.4: Multiple Testing}

S.4 applies the Benjamini-Hochberg procedure to the L-track's 1,050-test
battery (7 metrics $\times$ 50 states $\times$ 3 cycles).
At $q = 0.10$, BH correction reduces significant findings from 147
(uncorrected) to 89.
All clear gerrymanders (NC, WI, TX enacted) remain significant.
The 58 demoted cases were at the margin of significance under the
uncorrected test, suggesting they are likely false positives.

\subsection{Integration}

The three components integrate as follows in a complete analysis:

\begin{enumerate}
\item Generate ensemble with minimum size from S.3's 50-state table
\item Certify convergence using G.4's $\hat{R}$ and ESS diagnostics
\item Compute calibrated p-values using S.1's ESS-corrected formula
\item Apply S.4's BH procedure to the full test battery
\item Report BDS from S.2 as a court-facing summary
\end{enumerate}

This pipeline converts the current practice (``plan at Xth percentile'')
into litigation-grade inference (``plan is statistically significant at
FDR $q = 0.10$, with 80\% power to detect the observed effect size,
BDS = 0.97'').
